\subsection{Olympus ensemble}

Results for the standard DSI-ESMDA are presented in \cref{tab:olympus-dsi-esmda},
while results for the proposed Gated DSI-ESMDA are shown in \cref{tab:olympus-gated-dsi-esmda}.

Compared to the Egg case, priors of the Olympus ensemble presented extremely high data mismatches
(to the point of producing expurious deviations), 
which is a testment to the large spread in range this ensemble has.

\Cref{tab:olympus-dsi-esmda} confirms that localization had a pronounced positive effect on posterior coverage,
both in the history and validation periods, albeit at the risk of incurring high data mismatches.
Averaging across groundtruths,
localization increased the posterior coverage from 
40.00\% ($\pm$ 5.01\%) to 64.80\% ($\pm$ 5.75\%) in the history period,
and from
47.5\% ($\pm$ 18.53\%) to 87.8\% ($\pm$ 4.92\%) in the validation period,
while the average data mismatch increased from 
0.897 ($\pm$ 0.257) to 3.340 ($\pm$ 5.267).
%
Localization was also effective in decreasing the posterior MSE,
even managing to fair biased MSEs (groundtruths 15, 45).

\begin{table}[h!]
    \centering
    \IBGEtab{%
        \caption{%
            DSI-ESMDA results for the Olympus ensemble.
        }
        \label{tab:olympus-dsi-esmda}
    }{%
        \def\arraystretch{0.9}
        \begin{tabular}{ccccccc}
\toprule
\multirow{2}{*}{
    Gt.
}
&
\multirow{2}{*}{Loc.} 
& 
\multicolumn{2}{c}{OC} & 
\multirow{2}{*}{
    MSE
} & 
\multicolumn{2}{c}{Data mismatch} 
\\
& & His. & Val. &  & Avg. & Std. Dev.
\\
\midrule
\multirow{2}{*}{5} & No & 0.28 & 0.26 & 2.70 & 1.42 & 0.01 \\
 & Yes & 0.66 & 0.76 & 2.29 & 17.80 & 18.21 \\
\midrule
\multirow{2}{*}{10} & No & 0.40 & 0.59 & 1.83 & 0.98 & 0.02 \\
 & Yes & 0.59 & 0.85 & 1.56 & 1.30 & 0.29 \\
\midrule
\multirow{2}{*}{15} & No & 0.41 & 0.26 & 2.05 & 1.10 & 0.02 \\
 & Yes & 0.70 & 0.89 & 1.75 & 1.49 & 0.33 \\
\midrule
\multirow{2}{*}{20} & No & 0.42 & 0.63 & 1.11 & 0.62 & 0.02 \\
 & Yes & 0.57 & 0.86 & 1.09 & 0.72 & 0.06 \\
\midrule
\multirow{2}{*}{25} & No & 0.36 & 0.27 & 1.64 & 0.89 & 0.02 \\
 & Yes & 0.71 & 0.91 & 1.41 & 5.49 & 4.78 \\
\midrule
\multirow{2}{*}{30} & No & 0.43 & 0.78 & 1.15 & 0.63 & 0.01 \\
 & Yes & 0.60 & 0.87 & 1.10 & 0.75 & 0.09 \\
\midrule
\multirow{2}{*}{35} & No & 0.42 & 0.45 & 1.49 & 0.81 & 0.01 \\
 & Yes & 0.59 & 0.92 & 1.34 & 1.05 & 0.18 \\
\midrule
\multirow{2}{*}{40} & No & 0.45 & 0.56 & 1.41 & 0.77 & 0.01 \\
 & Yes & 0.72 & 0.88 & 1.24 & 1.76 & 0.74 \\
\midrule
\multirow{2}{*}{45} & No & 0.39 & 0.34 & 2.07 & 1.10 & 0.01 \\
 & Yes & 0.64 & 0.93 & 1.82 & 1.99 & 0.73 \\
\midrule
\multirow{2}{*}{50} & No & 0.45 & 0.61 & 1.18 & 0.65 & 0.01 \\
 & Yes & 0.70 & 0.91 & 1.12 & 1.05 & 0.23 \\
\bottomrule
\end{tabular}

    }{%
        \fonte{Self-elaboration.}
    }
\end{table}
%
\begin{table}[htbp]
    \centering
    \IBGEtab{%
        \caption{%
            Gated DSI-ESMDA ($\lambda=2$) results for the Olympus ensemble.
        }
        \label{tab:olympus-gated-dsi-esmda}
    }{%
        \def\arraystretch{0.9}
        \begin{tabular}{ccccccc}
\toprule
\multirow{2}{*}{
    Gt.
}
&
\multirow{2}{*}{Loc.} 
& 
\multicolumn{2}{c}{OC} & 
\multirow{2}{*}{
    MSE
} & 
\multicolumn{2}{c}{Data mismatch} 
\\
& & His. & Val. &  & Avg. & Std. dev.
\\
\midrule
\multirow{2}{*}{5} & No & 0.16 & 0.14 & 2.71 & 1.37 & 5.7 $\times 10^{-3}$ \\
 & Yes & 0.66 & 0.76 & 2.26 & 17.80 & 18.21 \\
\midrule
\multirow{2}{*}{10} & No & 0.29 & 0.47 & 1.84 & 0.96 & 0.01 \\
 & Yes & 0.57 & 0.87 & 1.62 & 1.31 & 0.30 \\
\midrule
\multirow{2}{*}{15} & No & 0.30 & 0.21 & 2.07 & 1.07 & 0.01 \\
 & Yes & 0.67 & 0.89 & 1.75 & 1.48 & 0.34 \\
\midrule
\multirow{2}{*}{20} & No & 0.32 & 0.55 & 1.14 & 0.61 & 0.01 \\
 & Yes & 0.54 & 0.86 & 1.15 & 0.74 & 0.04 \\
\midrule
\multirow{2}{*}{25} & No & 0.25 & 0.18 & 1.65 & 0.84 & 6.2 $\times 10^{-3}$ \\
 & Yes & 0.69 & 0.91 & 1.39 & 5.45 & 4.76 \\
\midrule
\multirow{2}{*}{30} & No & 0.34 & 0.71 & 1.17 & 0.63 & 0.01 \\
 & Yes & 0.55 & 0.91 & 1.14 & 0.77 & 0.08 \\
\midrule
\multirow{2}{*}{35} & No & 0.30 & 0.41 & 1.52 & 0.79 & 0.01 \\
 & Yes & 0.53 & 0.92 & 1.40 & 1.06 & 0.18 \\
\midrule
\multirow{2}{*}{40} & No & 0.29 & 0.34 & 1.41 & 0.73 & 7.8 $\times 10^{-3}$ \\
 & Yes & 0.70 & 0.89 & 1.25 & 1.74 & 0.72 \\
\midrule
\multirow{2}{*}{45} & No & 0.22 & 0.26 & 2.08 & 1.06 & 8.5 $\times 10^{-3}$ \\
 & Yes & 0.59 & 0.93 & 1.83 & 1.97 & 0.72 \\
\midrule
\multirow{2}{*}{50} & No & 0.33 & 0.47 & 1.19 & 0.63 & 7.8 $\times 10^{-3}$ \\
 & Yes & 0.67 & 0.91 & 1.13 & 1.05 & 0.23 \\
\bottomrule
\end{tabular}

    }{%
        \fonte{Self-elaboration.}
    }
\end{table}

\Cref{tab:olympus-gated-dsi-esmda} shows that the trend is similar for the proposed approach.
Averaging across groundtruths,
localization increased the posterior coverage from 
28.90\% ($\pm$ 6.54\%) to 61.7\% ($\pm$ 6.72\%) in the history period,
and from
37.40\% ($\pm$ 18.17\%) to 88.50\% ($\pm$ 4.90\%) in the validation period,
while the average data mismatch increased from 
0.869 ($\pm$ 0.246) to 3.337 ($\pm$ 5.263).
%
As with the standard DSI-ESMDA, localization was also effective in decreasing the posterior MSE,
managing to fair the same biased MSEs (groundtruths 15, 45).

% For this case study, the proposed gating was, at best, only as effective as the standard DSI-ESMDA in terms of posterior coverage and data mismatch.
% This could be due to poor tuning of the gating intensity hyperparameter, and remains a topic of future research.

\Cref{fig:olympus-summaries} summarizes results for the Olympus ensemble by averaging posterior coverage and MSE across groundtruths
for different gating intensities, with null gating corresponding to standard DSI-ESMDA.

As with the Egg case, intensifying the gating tends to increase the posterior coverage while trading off MSE.
Unlike the Egg case, however, MSE values of gated DSI-ESMDA are on average consistently higher than those of standard DSI-ESMDA.
Furthermore, validation coverage varied wildly, with localization greatly boosting the metric for null or mild gating.

\begin{figure}[htbp]
    \caption{%
        Summary of results for the Olympus ensemble.
        Posterior coverage and MSE are averaged across groundtruths.
    }
    \label{fig:olympus-summaries}
    \begin{subfigure}{.49\linewidth}
        \centering
        \includegraphics[width=\linewidth]{results/olympus/summary_plots/summary_history_coverage.pdf}
        \caption{Posterior coverage.}
    \end{subfigure}
    %
    \begin{subfigure}{.49\linewidth}
        \centering
        \includegraphics[width=\linewidth]{results/olympus/summary_plots/summary_validation_coverage.pdf}
        \caption{Validation coverage.}
    \end{subfigure}
    %
    \begin{subfigure}{.49\linewidth}
        \centering
        \includegraphics[width=\linewidth]{results/olympus/summary_plots/summary_history_mse.pdf}
        \caption{MSE.}
    \end{subfigure}
    \fonte{Self-elaboration.}
\end{figure}

The larger deviations and ranges on the average MSE and coverages are likely due, 
at least in some part, to the sampling of groundtruth realizations: FOPT-wise, half of them come from the tails of the distribution, 
which are weakly represented in the prior ensemble.
In order to explore the effects of gating on a more typical groundtruth,
Appendix C shows the prior and posterior distributions for realization 27, which is the ensemble's median FOPT.
The posterior is reasonably well-matched to the groundtruth.